Este trabajo abordó el Capacitated Vehicle Routing Problem, problema
$\mathcal{NP}$-hard de optimización combinatoria que consiste en diseñar
rutas de costo mínimo para una flota homogénea de vehículos con capacidad
limitada, exclusivamente mediante métodos exactos de programación matemática.

Se implementó el esquema clásico de Laporte y Nobert (1983) sobre la
formulación simétrica de dos índices, con generación iterativa de
desigualdades de capacidad redondeada y resolución de cada modelo entero
mediante \textit{Branch and Bound} a través del solver COIN-OR CBC. El método
certificó el óptimo global en cuatro de las cinco instancias evaluadas de
CVRPLIB (gap 0\,\%: \texttt{E-n13-k4}, \texttt{E-n22-k4}, \texttt{E-n23-k3} y
\texttt{A-n32-k5}), y en la quinta (\texttt{A-n33-k5}) obtuvo una solución
factible a 0{,}91\,\% del óptimo sin alcanzar el certificado dentro del
presupuesto de tiempo.

Los hallazgos confirman la doble naturaleza de los métodos exactos clásicos:
por una parte, proporcionan garantías de optimalidad verificables ---la
coincidencia exacta con las soluciones de referencia--- que ningún método
aproximado ofrece; por otra, su esfuerzo computacional crece de forma abrupta
con el tamaño de la instancia, al punto de que un solo nodo adicional desplazó
el problema fuera del rango resoluble. La implicación general es que el
esquema clásico resulta apropiado como herramienta de certificación sobre
instancias pequeñas y medianas, y como base conceptual ---relajación, cotas,
separación y ramificación--- de los métodos modernos que extienden su alcance,
mientras que las instancias mayores exigen los enfoques \textit{Branch and
Cut} y \textit{Branch, Cut and Price} revisados en el estado del arte, o bien
métodos aproximados cuya exploración corresponde a la siguiente etapa de este
trabajo.